\section{Conclusion} \label{conclu}

Instance normalization is best understood as a conditional experimental choice,
not a universal forecasting default. It removes look-back offset and scale,
which can improve robustness when those quantities vary across dates or series
without changing the normalized forecasting relationship. The same operation
can be detrimental when absolute level and scale are predictive. Normalized
backpropagation adds a separate effect by balancing gradient contributions from
series with different scales, while RevIN's symmetric affine parameters must be
judged independently.

Our expanded protocol tests these mechanisms across datasets, horizons,
backbones, future dates, unseen series, and random seeds. Reporting seed
uncertainty and within-run error variance is essential because small average
gains may otherwise be indistinguishable from training variability. Finally,
the standard-versus-RevIN oracle measures the headroom available from
choosing normalization per setting, while remaining explicitly optimistic.
Together, these analyses aim to replace the blanket question ``does RevIN
work?'' with the more useful question ``under which data conditions does
instance normalization help?''
